\documentclass[11pt,a4paper]{article}

\usepackage[utf8]{inputenc}
\usepackage[T1]{fontenc}
\usepackage{amsmath,amssymb,amsthm}
\usepackage{bm}
\usepackage{geometry}
\geometry{margin=2.5cm}
\usepackage{braket}
\usepackage{hyperref}
\usepackage{xcolor}
\usepackage{listings}

\lstset{
  basicstyle=\ttfamily\small,
  breaklines=true,
  columns=fullflexible,
  frame=single,
  keywordstyle=\color{blue!60!black}\bfseries,
  commentstyle=\color{green!50!black}\itshape,
  language=Python,
}

% Common shorthands
\newcommand{\rr}{\mathbf{r}}
\newcommand{\qq}{\mathbf{q}}
\newcommand{\GG}{\mathbf{G}}
\newcommand{\kk}{\mathbf{k}}
\newcommand{\Gpar}{\mathbf{G}_{\!\parallel}}
\newcommand{\qpar}{\mathbf{q}_{\parallel}}
\newcommand{\Gz}{G_{z}}
\newcommand{\dd}{\mathrm{d}}
\newcommand{\im}{\mathrm{i}}
\newcommand{\half}{\tfrac{1}{2}}
\newcommand{\Tr}{\mathrm{Tr}}
\newcommand{\Heter}{\text{HS}}
\newcommand{\sgn}{\mathrm{sgn}}
\newcommand{\eps}{\varepsilon}
\newcommand{\eq}[1]{Eq.~(\ref{eq:#1})}

\title{The $G\Delta W$ method in the mQEH basis: \\ mathematical analysis of the \texttt{GWmQEHCorrection} implementation}
\author{Branch \texttt{claude/implement-gdeltaw-mqeh-E3yyU} of GPAW}
\date{}

\begin{document}
\maketitle

\begin{abstract}
We document the mathematics behind \texttt{gpaw.response.gwqeh.GWmQEHCorrection},
which extends the $G\Delta W$ method of Winther and
Thygesen [2D~Mater.~\textbf{4}, 025059 (2017)] from the monopole approximation
to the full multi-pole quantum-electrostatic heterostructure (mQEH)
basis.  The original $G\Delta W$ method assumes that the change in the
screened Coulomb interaction $\Delta W^i$ of a 2D layer~$i$ embedded in
a van der Waals heterostructure is well described by its in-plane
macroscopic ($\GG=\GG'=0$) component, equivalent to keeping a single
monopole density basis function per layer.  In the mQEH generalisation
the building block of each layer carries multiple density and
potential basis functions $\{\rho_{\alpha},\phi_{\alpha}\}$, and
$\Delta W^i$ is a matrix in this basis.  We derive how this matrix
enters the GW self-energy when the in-plane plane-wave expansion is
truncated by a cutoff $E_\text{cut}^{\text{mQEH}}$, give explicit
expressions for the expansion coefficients of the GPAW pair densities,
and map every term to the lines of the implementation.
\end{abstract}

\section{Recap of the $G\Delta W$ method}

The screened Coulomb interaction
\begin{equation}
W(\rr,\rr',\omega) = \int \dd\rr''\, \eps^{-1}(\rr,\rr'',\omega)\,
   V(\rr'',\rr'),
\label{eq:W-def}
\end{equation}
splits in a heterostructure of 2D layers as
\begin{equation}
W^{i,\Heter}(\rr,\rr',\omega) =
   W^i(\rr,\rr',\omega) + \Delta W^i(\rr,\rr',\omega),
\label{eq:dW-decomp}
\end{equation}
where $W^i$ is the freestanding-monolayer screened interaction and
$\Delta W^i$ is the change in screening due to the surrounding layers
(Eq.~(5) of the paper).  The QP correction due to the environment is
\begin{equation}
\Delta E_{n\kk}^{i,\text{scr}} = Z_{n\kk}^i\,
   \mathrm{Re}\,\Delta\Sigma_{n\kk}^i(\omega=\eps_{n\kk}),
\label{eq:dE-Z-dSigma}
\end{equation}
with $Z_{n\kk}^i = (1 - \partial_\omega\Sigma_{n\kk})^{-1}$ the GW
renormalisation factor (Eq.~(6) of the paper).

In a plane-wave representation the full GW self-energy on top of
DFT eigenstates reads
\begin{equation}
\Sigma_{n\kk}=\frac{1}{(2\pi)^3}\!\int_\text{BZ}\!\dd\qq\,
\frac{\im}{2\pi}\!\int\dd\omega'\sum_{m,\GG\GG'}
\frac{\rho^{m,\kk-\qq}_{n,\kk}(\GG)\,
      \bigl[\rho^{m,\kk-\qq}_{n,\kk}(\GG')\bigr]^{*}}
     {\omega+\omega'-\eps_{m\kk-\qq}-\im\eta\,\sgn(\eps_{m\kk-\qq}-\mu)}
\,\widetilde W_{\GG\GG'}(\qq,\omega'),
\label{eq:Sigma-PW}
\end{equation}
with $\widetilde W = W - V$ (Eq.~(9) of the paper) and the pair density
\begin{equation}
\rho^{m,\kk-\qq}_{n,\kk}(\GG) =
   \braket{n\kk | \, e^{\im (\qq+\GG)\!\cdot\!\rr}\,| m,\kk-\qq}.
\label{eq:pair-density}
\end{equation}
The implementation actually computes the conjugate convention
$\rho_{n,\kk}^{m,\kk-\qq}(\GG)$ via
\texttt{ActualPairDensityCalculator.calculate\_pair\_density}, which we
will use below.

\paragraph{Original (monopole) approximation.}
$\Delta W^i$ is generated by charges \emph{outside} layer~$i$ and
therefore varies slowly across the (thin) layer.  Winther \& Thygesen
exploit this to retain only the macroscopic component
$\GG=\GG'=0$ in \eq{Sigma-PW}, which is the average of
$\Delta W^i(z,z',\qq,\omega)$ over the layer thickness $d_i$:
\begin{equation}
\Delta W^i(\qq,\omega) =
\frac{1}{d_i^2}\int_{z_i-d_i/2}^{z_i+d_i/2}\!\!\!\!\int
W^{\Heter}(z,z',\qq,\omega) - W^i(z,z',\qq,\omega)\,\dd z\,\dd z'.
\label{eq:dW-monopole}
\end{equation}
The resulting self-energy correction is the central equation of the
original $G\Delta W$ paper (its Eq.~(12)),
\begin{equation}
\Delta\Sigma_{n\kk}^{i,\text{mono}}\;=\;
\frac{1}{(2\pi)^3}\!\int_\text{BZ}\!\!\dd\qq\,
\frac{\im}{2\pi}\!\int\dd\omega'\sum_{m}
\frac{|\rho^{m,\kk-\qq}_{n,\kk}(\GG=0)|^{2}}
     {\omega+\omega'-\eps_{m\kk-\qq}-\im\eta\,\sgn(\eps_{m\kk-\qq}-\mu)}
\,\Delta W^i(\qq,\omega').
\label{eq:Sigma-monopole}
\end{equation}
This is the formula implemented in the parent class
\texttt{GWQEHCorrection.calculate\_QEH} of
\texttt{gpaw/response/gwqeh.py}: the descriptor is built with
\texttt{self.ecut = 0.1\,Ha} so that the plane-wave basis contains
exactly the single vector $\GG=0$.

\section{The mQEH basis}

The QEH model represents the linear density response of an isolated
2D layer~$i$ in a small basis of \emph{density basis functions}
$\{\rho^{i}_{\alpha}(z;q_\parallel)\}_{\alpha=1}^{N_i}$, which depend
on the magnitude of the in-plane momentum
$q_\parallel = |\qpar|$.  Each density basis function generates a
\emph{potential basis function} $\phi^{i}_{\alpha}(z;q_\parallel)$
through the 1-D Poisson equation in the in-plane Fourier
representation,
\begin{equation}
\bigl(-\partial_{z}^{2}+q_\parallel^{2}\bigr)\,
\phi^{i}_{\alpha}(z;q_\parallel)
\;=\; 4\pi\, \rho^{i}_{\alpha}(z;q_\parallel).
\label{eq:Poisson}
\end{equation}
The two sets are biorthogonal,
\begin{equation}
\int \dd z\;\phi^{i*}_{\alpha}(z;q_\parallel)\,
\rho^{i}_{\beta}(z;q_\parallel)\;=\;\delta_{\alpha\beta},
\label{eq:biorth}
\end{equation}
a normalisation that is enforced inside the QEH building-block
construction.  The original monopole QEH corresponds to $N_i=1$ with
$\rho^{i}_{1}$ the constant in-plane charge profile of the layer.

In this basis the response function and the screened interaction
become \emph{matrices}.  For a full heterostructure with layers
$l=1,\ldots,L$ and total basis dimension
$N=\sum_{l} N_l$ the building blocks couple as
\begin{equation}
\hat W^{\Heter}(q_\parallel,\omega) =
\hat V(q_\parallel)+\hat V(q_\parallel)\,\hat\chi^{\Heter}(q_\parallel,\omega)\,
\hat V(q_\parallel),
\label{eq:W-matrix}
\end{equation}
where carets denote $N\times N$ matrices.  The change in screening
acting on layer~$i$ is the on-layer block of
\begin{equation}
\Delta W^{i}_{\alpha\beta}(q_\parallel,\omega)\;=\;
W^{\Heter}_{(i\alpha)(i\beta)}(q_\parallel,\omega)
\;-\;W^{i,0}_{\alpha\beta}(q_\parallel,\omega),
\qquad \alpha,\beta=1,\ldots,N_i,
\label{eq:dW-matrix}
\end{equation}
where $W^{i,0}$ is computed for the isolated layer.  This is the
quantity produced by the QEH code through
\texttt{QEH.heterostructure(\ldots).get\_screened\_potential()} and is
sliced in \texttt{calculate\_W\_QEH} of
\texttt{GWmQEHCorrection} (\texttt{gwqeh.py} lines~792--813).

The corresponding decomposition of $\Delta W^i$ in mixed
real-space/in-plane-momentum representation is the one assumed
throughout this work:
\begin{equation}
\Delta W^{i}(\qpar; z, z'; \omega)
\;\approx\;\sum_{\alpha\beta=1}^{N_i}
\phi^{i}_{\alpha}(z;q_\parallel)\,
\Delta W^{i}_{\alpha\beta}(q_\parallel,\omega)\,
\phi^{i*}_{\beta}(z';q_\parallel).
\label{eq:dW-zz-expansion}
\end{equation}
The monopole limit is recovered by retaining only
$\alpha=\beta=1$ with $\phi^{i}_{1}\equiv 1/d_i$ (constant inside the
layer, zero outside), which reduces \eq{dW-zz-expansion} to the
spatial average~(\ref{eq:dW-monopole}).

\section{Self-energy in the mQEH basis}

\subsection{Plane-wave to mixed representation}

Write the plane-wave vector $\GG=(\Gpar,\Gz)$ and the
in-plane reciprocal vector $\qpar+\Gpar$.  In a slab supercell of
height $L_z$ and in-plane area $A$ (so $\Omega = A L_z$), the
plane-wave matrix element of $\Delta W^i$ is, using
\eq{dW-zz-expansion} and assuming full in-plane translational
invariance of the heterostructure ($\Delta W^i$ depending only on the
in-plane separation, hence diagonal in $\Gpar$),
\begin{equation}
\Delta W^i_{\GG\GG'}(\qq,\omega) \;=\;
\frac{\delta_{\Gpar,\Gpar'}}{L_z}
\!\sum_{\alpha\beta}
\tilde\phi^{i}_{\alpha}(\Gz;\,|\qpar+\Gpar|)\,
\Delta W^{i}_{\alpha\beta}(|\qpar+\Gpar|,\omega)\,
\bigl[\tilde\phi^{i}_{\beta}(\Gz';\,|\qpar+\Gpar|)\bigr]^{\!*},
\label{eq:dWGG}
\end{equation}
with
\begin{equation}
\tilde\phi^{i}_{\alpha}(\Gz;q_\parallel)
\;=\;\int_{0}^{L_z}\!\dd z\, e^{-\im\Gz z}\,
\phi^{i}_{\alpha}(z;q_\parallel).
\label{eq:phi-FT}
\end{equation}
This is the key identity that makes the mQEH basis tractable:
$\Delta W^i$ becomes a small ($N_i\!\times\!N_i$) matrix in
$(\alpha,\beta)$, and the dependence on $\Gz,\Gz'$ is absorbed into
the universal layer functions $\phi^{i}_{\alpha}$.

\subsection{Projection of the pair densities}

Inserting \eq{dWGG} into \eq{Sigma-PW} and changing the order of
summations gives
\begin{equation}
\Delta\Sigma^{i,\text{mQEH}}_{n\kk} \;=\;
\frac{1}{(2\pi)^3 L_z}\!\!
\int_\text{BZ}\!\dd\qq\,\frac{\im}{2\pi}\!\!\int\dd\omega'\!
\sum_{m,\,\Gpar}\,
\frac{1}{\omega+\omega'-\eps_{m\kk-\qq}-\im\eta\,\sgn(\dots)}
\sum_{\alpha\beta}
\mathcal{C}^{(nm)*}_{\alpha}(\qpar+\Gpar)\,
\Delta W^i_{\alpha\beta}\bigl(|\qpar+\Gpar|,\omega'\bigr)\,
\mathcal{C}^{(nm)}_{\beta}(\qpar+\Gpar),
\label{eq:Sigma-mQEH-final}
\end{equation}
with the projected pair-density coefficients
\begin{equation}
\boxed{\;
\mathcal{C}^{(nm)}_{\alpha}(\qpar+\Gpar) \;\equiv\;
\sum_{\Gz}\,
\rho^{m,\kk-\qq}_{n,\kk}(\Gpar,\Gz)\;
\tilde\phi^{i}_{\alpha}(\Gz;|\qpar+\Gpar|).
\;}
\label{eq:C-def}
\end{equation}

Using Parseval's identity in the out-of-plane direction,
\eq{C-def} is equivalent to a real-space overlap with the potential
basis function:
\begin{equation}
\mathcal{C}^{(nm)}_{\alpha}(\qpar+\Gpar) \;=\;
\int_{0}^{L_z}\dd z\;
\phi^{i}_{\alpha}(z;|\qpar+\Gpar|)\,
\bar\rho^{m,\kk-\qq}_{n,\kk}(\Gpar;z),
\label{eq:C-realspace}
\end{equation}
where the in-plane Fourier-component of the pair density, resolved in
$z$, is
\begin{equation}
\bar\rho^{m,\kk-\qq}_{n,\kk}(\Gpar;z) \;=\;
\frac{1}{L_z}\sum_{\Gz}
\rho^{m,\kk-\qq}_{n,\kk}(\Gpar,\Gz)\, e^{+\im\Gz z}.
\label{eq:rho-z}
\end{equation}
\eq{C-realspace} is the form used in the implementation, because the
mQEH library returns the basis functions $\phi^{i}_{\alpha}(z)$
sampled on a real-space $z$-grid common to the whole heterostructure.

\subsection{The $1/A$ vs $1/\Omega$ factor}

The original PW formula \eq{Sigma-PW} carries a $1/\Omega$ prefactor
through the BZ measure $\dd\qq/(2\pi)^3$ together with the unit
3-D-volume normalisation of $|n\kk\rangle$.  After integrating the
$\Gz$ sums into $\bar\rho(\Gpar;z)$, the leftover sum over the
$\Gpar$ grid is the in-plane PW expansion and contributes the
2-D BZ measure with prefactor $1/A$.  This is the origin of the
explicit $1/L_z$ in \eq{Sigma-mQEH-final} and of the multiplication
by $L = L_z$ of the mQEH matrix in the code
(\texttt{gwqeh.py} line~1083):
\begin{lstlisting}[language=Python]
dW_wab = self._eval_dW_on_gwgrid(q_abs)
dW_wab *= L
\end{lstlisting}
After this rescaling the original prefactor
\texttt{x = 1/(N\_q\,2\,pi\,$\Omega$)} (line~1258, identical to the
parent class) can be reused unchanged.

\subsection{Frequency representation}

As in the parent class the temporal Fourier integral in
\eq{Sigma-mQEH-final} is evaluated through a Hilbert transform.  Let
$\hat\omega = \omega+\omega'$ and use the standard decomposition
$\Delta W^i_{\alpha\beta}(q_\parallel,\omega') = \Delta
W^{+}_{\alpha\beta}+\Delta W^{-}_{\alpha\beta}$ corresponding to the
two $\pm\im\eta$ branches.  Each branch is then convolved via the
non-uniform spectral grid $\{\omega_w\}$.  Concretely the code stores
both transforms in a single tensor
\texttt{Wpm\_Gpar} of shape $(n_{\Gpar},2N_\omega,N_i,N_i)$:
\begin{lstlisting}[language=Python]
Wpm_Gpar[ig, :nw] = dW_wab
Wpm_Gpar[ig, nw:] = dW_wab.copy()
self.htp(Wpm_Gpar[ig, :nw])
self.htm(Wpm_Gpar[ig, nw:])
\end{lstlisting}
After this transformation the contribution from each electronic state
$m$ reduces to a linear interpolation between two adjacent frequency
points (\texttt{w} and \texttt{w+1}), exactly as in the standard
GW code, but with the scalar $W_w$ replaced by an
$N_i\!\times\!N_i$ matrix $W^{\pm}_{w,\alpha\beta}$.

\section{Mapping the math to \texttt{GWmQEHCorrection}}

We list the correspondence between each step in
\texttt{calculate\_QEH} / \texttt{\_calculate\_sigma\_mqeh}
(\texttt{gpaw/response/gwqeh.py}) and the equations above.

\paragraph{(M1) Building the in-plane PW basis.}
The cutoff is raised from the monopole limit
$E_\text{cut}=0.1\,\text{Ha}$ to a user-controlled value
$E_\text{cut}^{\text{mQEH}}$, exposed as
\texttt{ecut\_mqeh} (default $50\,\text{eV}$, line~677):
\begin{lstlisting}[language=Python]
pd0 = SingleQPWDescriptor.from_q(q_c, self.ecut_mqeh,
                                 self.gs.gd, gammacentered=True)
G_Gv  = pd0.get_reciprocal_vectors()           # q + G
G0_Gv = pd0.get_reciprocal_vectors(add_q=False)
\end{lstlisting}
The vectors are partitioned by their in-plane component
$\Gpar=(G_x,G_y)$:
\begin{lstlisting}[language=Python]
Gpar_rounded = np.round(G0_Gv[:, :2], decimals=8)
unique_Gpar, inverse_idx = np.unique(
    Gpar_rounded, axis=0, return_inverse=True)
\end{lstlisting}
This realises the partitioning of the inner sum
$\sum_{\GG\GG'}\rightarrow \sum_{\Gpar}\sum_{\Gz\Gz'}$ in
\eq{Sigma-PW} that leads to \eq{Sigma-mQEH-final}.  An assertion on
the cell metric (lines~718--722) ensures that the third lattice
vector is along Cartesian $\hat z$, so that this Cartesian split
coincides with the lattice partition.

\paragraph{(M2) Evaluating $\Delta W^i_{\alpha\beta}$ at
$|\qpar+\Gpar|$.}  The mQEH matrix data are pre-interpolated by
\texttt{\_interpolate\_mqeh\_data} (lines~871--964): one cubic spline
in $\omega$ moves the data onto the GW frequency grid; then a cubic
spline in $q_\parallel$ over the resulting array lets us evaluate
$\Delta W^i_{\alpha\beta}(|\qpar+\Gpar|,\omega)$ for every member of
the $\Gpar$-grid with a single \texttt{CubicSpline} call.  The result
is the rescaled matrix
$L_z\,\Delta W^i_{\alpha\beta}$ stored in \texttt{Wpm\_Gpar}, ready
for the Hilbert transform.

The small-$q$ behaviour is handled exactly as in the parent class
(\texttt{get\_W\_on\_grid}, lines~933--964): inside a disk of radius
\texttt{q\_cut} around $\qpar+\Gpar=0$ the spline is replaced by a
weighted average over the sampled small-$q$ points
\begin{equation}
\overline{\Delta W}^{\,q\to 0}_{\alpha\beta}(\omega) =
\sum_{q_0\le q_\text{cut}} w(q_0)\, \Delta W^i_{\alpha\beta}(q_0,\omega),
\qquad
w(q_0) \propto q_0,
\label{eq:q0-average}
\end{equation}
mirroring the integration of $W(q)$ over the $\qq=0$ disk in the
parent class.  The same logic is now applied to every matrix element
of $\Delta W^i_{\alpha\beta}$.

\paragraph{(M3) Computing $\bar\rho(\Gpar;z)$.}  The pair density
\eq{pair-density} is computed by
\texttt{ActualPairDensityCalculator} in the same call that the parent
class uses, but now with the larger PW basis.  For each unique
$\Gpar$, the implementation collects the indices of the $\Gz$ members
and applies \eq{rho-z}:
\begin{lstlisting}[language=Python]
phase_zg = exp(1j * np.outer(z_qeh, Gz_values))   # exp(+i Gz z)
n_m_gz   = n_mG[:, G_indices]                     # rho(Gpar, Gz)
rho_mz   = n_m_gz @ phase_zg.T / Lz               # bar-rho(Gpar; z)
\end{lstlisting}
Note the $1/L_z$ in line~1274, exactly the factor in \eq{rho-z}; the
phase $e^{+\im\Gz z}$ uses the sign convention of
\texttt{calculate\_pair\_density}.  The $z$-grid \texttt{z\_qeh} and
the spacing \texttt{dz\_qeh} are imported from the QEH heterostructure
object together with the basis functions, so the same discretisation
is used for $\rho$ and $\phi$.

The phase factors $\exp(\im\Gz z)$ depend only on the $\Gpar$ group
and the q-point geometry, not on band index or symmetry, so they are
precomputed once per IBZ q-point and reused inside
\texttt{\_calculate\_sigma\_mqeh}.  This is the optimisation
introduced in commit \texttt{fee44a26}, \emph{Vectorize mQEH
interpolation; hoist z-Fourier phase out of band loop}.

\paragraph{(M4) Projection coefficients $\mathcal{C}^{(nm)}_{\alpha}$.}
\eq{C-realspace} is implemented as a single matrix product:
\begin{lstlisting}[language=Python]
phi_az = phi_Gpar_az[ig]                                # (N_i, n_z)
C_mGpar_a[:, ig, :] = rho_mz @ phi_az.conj().T * self.dz_qeh
\end{lstlisting}
The conjugation on \texttt{phi\_az} implements the $\phi^{i}_{\alpha}$
factor in \eq{C-realspace} (the basis functions returned by QEH are
defined to be real for real $q_\parallel$, but the code keeps the
complex form for generality).  The factor \texttt{dz\_qeh} is the
trapezoidal-rule weight that turns the sum over the QEH $z$-grid into
$\int dz$.

\paragraph{(M5) Quadratic form and self-energy.}
The inner double sum in \eq{Sigma-mQEH-final},
\begin{equation}
\mathcal{T}_{nm}(\qpar+\Gpar,\omega') \;=\;
\sum_{\alpha\beta}
\mathcal{C}^{(nm)*}_{\alpha}\,
\Delta W^{i}_{\alpha\beta}\,
\mathcal{C}^{(nm)}_{\beta},
\label{eq:Tnm}
\end{equation}
is evaluated for each transition $m$ and each $\Gpar$ as a small
quadratic form,
\begin{lstlisting}[language=Python]
sigma1 += p * (C_a.conj() @ W1_ab @ C_a).imag
sigma2 += p * (C_a.conj() @ W2_ab @ C_a).imag
\end{lstlisting}
with the same Hilbert-transformed indexing \texttt{s*nw+w} used by the
parent class.  The two terms \texttt{sigma1}, \texttt{sigma2} are
linearly interpolated in $\omega$ to obtain $\Sigma_{n\kk}$ and
similarly for $\partial_\omega\Sigma_{n\kk}$.  Summation over
$\Gpar$ is the explicit Python loop \texttt{for ig in
range(n\_Gpar)}.

\paragraph{(M6) Compatibility with the parent class.}
The class inherits from \texttt{GWQEHCorrection} and overrides only
the three methods that change in the mQEH version:
\texttt{calculate\_W\_QEH} (returns the full matrix data),
\texttt{get\_W\_on\_grid} (interpolates the matrix too), and
\texttt{calculate\_QEH} (the self-energy loop with the projection
described above).  The monopole scalar $\Delta W^i_{00}$ is still
returned to the parent class so that the parent's restart logic and
zero-$q$ averaging continue to work; the additional matrix data is
serialised side-by-side under the same NPZ key
prefix ``\texttt{\_dW\_qw.npz}'' so that a restart reloads the full
problem (\texttt{\_try\_load\_mqeh\_npz}, lines~740--764).

\section{Reduction to the monopole limit}

To check the construction, consider $N_i=1$ with a constant potential
basis function $\phi^{i}_{1}(z)=1/d_i$ inside layer~$i$ and zero
outside.  The Fourier transform \eq{phi-FT} becomes
\begin{equation}
\tilde\phi^{i}_{1}(\Gz) = \frac{1}{d_i}\int_{z_i-d_i/2}^{z_i+d_i/2}
e^{-\im\Gz z}\dd z =
e^{-\im\Gz z_i}\,\mathrm{sinc}\!\bigl(\Gz d_i/2\bigr).
\label{eq:phi-tilde-mono}
\end{equation}
In the limit of a vanishingly thin layer this is unity for all $\Gz$
in the supercell, so
\begin{equation}
\mathcal{C}^{(nm)}_{1}(\qpar) \;=\;
\sum_{\Gz} \rho^{m,\kk-\qq}_{n,\kk}(\Gpar=0,\Gz)\cdot 1
\;\to\;
L_z\,\bar\rho^{m,\kk-\qq}_{n,\kk}(\Gpar=0;\,z\!\in\!\text{layer})
\;=\;\rho^{m,\kk-\qq}_{n,\kk}(\GG=0),
\end{equation}
i.e. the macroscopic pair density of \eq{Sigma-monopole}.  Together
with \eq{dW-zz-expansion} reducing to \eq{dW-monopole}, the mQEH
self-energy \eq{Sigma-mQEH-final} reduces to the parent-class
expression.  This is also tested numerically in
\texttt{test\_gwqeh.py} (\texttt{test\_mqeh\_zero\_dW},
\texttt{test\_mqeh\_linearity\_in\_dW},
\texttt{test\_mqeh\_physical\_signs}) using the synthetic data
generator \texttt{\_make\_synthetic\_mqeh\_data} with
\texttt{nbasis=1}.

\section{Summary of equations and implementation}

Table~\ref{tab:summary} collects the correspondence between the
mathematical objects and the implementation.

\begin{table}[h]
\centering
\small
\begin{tabular}{l l l}
\hline
Symbol & Meaning & Code object \\
\hline
$\rho^{i}_{\alpha}(z;q_\parallel)$ & density basis function of layer
$i$ & \texttt{drho\_qzi\_target}\\
$\phi^{i}_{\alpha}(z;q_\parallel)$ & potential basis function (dual
to $\rho^{i}_{\alpha}$) & \texttt{phi\_qiz\_target}\\
$\Delta W^{i}_{\alpha\beta}(q_\parallel,\omega)$ &
on-layer block of the mQEH $\Delta W$
& \texttt{dW\_qw\_matrix}\\
$\Gpar$ partition of $\GG$ & in-plane / out-of-plane split &
\texttt{unique\_Gpar}, \texttt{inverse\_idx}\\
$\bar\rho(\Gpar;z)$ & $z$-resolved pair density, \eq{rho-z} &
\texttt{rho\_mz}\\
$\mathcal{C}^{(nm)}_{\alpha}$ & projection on potential basis,
\eq{C-realspace} & \texttt{C\_mGpar\_a}\\
$L_z\,\Delta W^{i}_{\alpha\beta}$ & rescaled matrix used in
self-energy & \texttt{Wpm\_Gpar}\\
$\mathcal{T}_{nm}$ in \eq{Tnm} & matrix quadratic form &
\texttt{(C\_a.conj() @ W\_ab @ C\_a).imag}\\
$\Delta\Sigma^{i,\text{mQEH}}_{n\kk}$ & self-energy correction,
\eq{Sigma-mQEH-final} & \texttt{sigma\_sin}\\
\hline
\end{tabular}
\caption{Mathematical objects of \eq{Sigma-mQEH-final} and their
names in \texttt{gpaw/response/gwqeh.py}.}
\label{tab:summary}
\end{table}

\paragraph{Final self-energy formula.}  The expression actually
evaluated by \texttt{GWmQEHCorrection.calculate\_QEH} is
\begin{equation}
\boxed{\;
\Delta\Sigma^{i,\text{mQEH}}_{n\kk}\;=\;
\frac{1}{N_{\mathbf{q}}\,2\pi\,\Omega}
\sum_{\qq\in\text{IBZ}}\sum_{m,\,\Gpar}
\mathcal{I}_{nm}\!\bigl[
L_z\,\mathcal{C}^{(nm)*}_{\alpha}\,
\Delta W^{i}_{\alpha\beta}(|\qpar+\Gpar|)\,
\mathcal{C}^{(nm)}_{\beta}
\bigr](\eps_{n\kk}),
\;}
\label{eq:Sigma-code}
\end{equation}
where $\mathcal{I}_{nm}[\,\cdot\,](\omega)$ stands for the
Hilbert-transformed frequency convolution implemented by the
\texttt{HilbertTransform} class and the linear interpolation in
$\omega$ along the non-uniform grid \texttt{omega\_w}.  The
prefactor $1/(N_{\mathbf{q}}\,2\pi\,\Omega)$ is exactly the
\texttt{x} variable of \texttt{\_calculate\_sigma\_mqeh}; combined
with the $L_z$ already absorbed in $\Delta W$ it yields the
$1/(N_{\mathbf{q}}\,2\pi\,A)$ in-plane normalisation of
\eq{Sigma-mQEH-final}.

The QP correction is finally
\begin{equation}
\Delta E^{i,\text{scr,mQEH}}_{n\kk}
\;=\;Z_{n\kk}^{i}\,
\mathrm{Re}\,\Delta\Sigma^{i,\text{mQEH}}_{n\kk}(\omega=\eps_{n\kk}),
\end{equation}
identical in form to \eq{dE-Z-dSigma} of the original paper but with
$\Delta\Sigma$ now containing the full mQEH multi-pole response.  The
renormalisation factor $Z$ is taken from the freestanding-monolayer
$G_0W_0$ calculation, \texttt{gwfile}, exactly as in the parent class
(\texttt{calculate\_qp\_correction}, lines~379--408).

\section*{References}

\begin{description}
\item[\textbullet] K. T. Winther, K. S. Thygesen,
  \emph{Band structure engineering in van der Waals heterostructures
  via dielectric screening: the $G\Delta W$ method}, 2D Mater.\ 4,
  025059 (2017), \href{https://arxiv.org/abs/1703.03188}{arXiv:1703.03188}.
\item[\textbullet] K. Andersen, S. Latini, K. S. Thygesen,
  \emph{Dielectric genome of van der Waals heterostructures},
  Nano Lett.\ 15, 4616 (2015).
\item[\textbullet] Implementation:
  \texttt{gpaw/response/gwqeh.py}, class
  \texttt{GWmQEHCorrection},
  branch \texttt{claude/implement-gdeltaw-mqeh-E3yyU}.
\end{description}

\end{document}
